\documentclass[tikz,border=8pt]{standalone}
\usepackage[no-math]{fontspec}
\usepackage[brazil]{babel}
\usepackage{amsmath,amssymb}
\IfFileExists{C:/Windows/Fonts/Lato_Regular.ttf}{%
  \setmainfont{Lato}[
    Path = C:/Windows/Fonts/,
    UprightFont = Lato_Regular.ttf,
    BoldFont = Lato_Bold.ttf,
    ItalicFont = Lato_Italic.ttf,
    BoldItalicFont = Lato_BoldItalic.ttf
  ]
  \setsansfont{Lato}[
    Path = C:/Windows/Fonts/,
    UprightFont = Lato_Regular.ttf,
    BoldFont = Lato_Bold.ttf,
    ItalicFont = Lato_Italic.ttf,
    BoldItalicFont = Lato_BoldItalic.ttf
  ]
}{%
  \RequirePackage[default]{lato}
}
\renewcommand{\familydefault}{\sfdefault}

\usepackage{xcolor}
\usepackage{tikz}
\usepackage{pgfplots}
\pgfplotsset{compat=1.18}
\usetikzlibrary{arrows.meta,positioning,shapes.geometric,fit,backgrounds,calc,decorations.pathreplacing,patterns}

% Paleta institucional UEFS / CHF662
\definecolor{uefsNavy}{HTML}{1B2A4A}
\definecolor{uefsTeal}{HTML}{0E7C7B}
\definecolor{uefsForest}{HTML}{2D6A4F}
\definecolor{uefsGold}{HTML}{C4983B}
\definecolor{uefsGoldDark}{HTML}{8F6B1E}
\definecolor{uefsAlert}{HTML}{B02A37}
\definecolor{uefsOrange}{HTML}{D95F02}
\definecolor{uefsInk}{HTML}{22252A}
\definecolor{uefsMuted}{HTML}{5C6470}
\definecolor{uefsWater}{HTML}{2E86AB}
\definecolor{uefsSand}{HTML}{D9B382}
\definecolor{bgNavy}{HTML}{F2F5F9}
\definecolor{bgTeal}{HTML}{E9F5F4}
\definecolor{bgForest}{HTML}{EEF5F0}
\definecolor{bgGold}{HTML}{FCF7EB}
\definecolor{bgAlert}{HTML}{FDF2F2}
\definecolor{bgGray}{HTML}{F8FAFC}
\definecolor{lineGray}{HTML}{CFD8DC}

\tikzset{
  >=Stealth,
  aaCard/.style={draw=#1, fill=white, rounded corners=4pt, line width=1pt, inner sep=6pt, align=center, font=\small},
  aaBox/.style={draw=#1, rounded corners=4pt, line width=1.1pt, inner sep=6pt, align=center, font=\footnotesize\sffamily, text=uefsInk},
  aaTitle/.style={font=\normalsize\bfseries\sffamily, text=uefsNavy, align=center},
  aaSub/.style={font=\footnotesize\sffamily, text=uefsMuted, align=center},
  aaArrow/.style={->, >=Stealth, line width=1.4pt},
  aaDArrow/.style={<->, >=Stealth, line width=1.4pt}
}


\begin{document}
\begin{tikzpicture}[font=\sffamily]

% Título principal
\node[font=\large\bfseries\sffamily, text=uefsNavy, align=center] at (0, 3.65)
  {Produção do Espaço Geográfico: Meios Geográficos (Milton Santos, 1996)};

% =========================================================
% PAINEL 1: Meio Natural (Primeira Natureza)
% =========================================================
\begin{scope}[shift={(-5.8, 1.05)}]
  \draw[draw=uefsForest, fill=bgForest, rounded corners=6pt, line width=1.2pt] (-2.45, -1.75) rectangle (2.45, 1.85);
  \fill[uefsForest, rounded corners=5pt] (-2.45, 1.05) rectangle (2.45, 1.85);
  \fill[uefsForest] (-2.45, 1.05) rectangle (2.45, 1.35);
  \node[align=center, font=\small\bfseries\sffamily, text=white] at (0, 1.45)
    {1. Meio Natural\\[-1pt]\footnotesize\normalfont(Primeira Natureza)};

  \node[anchor=north west, text width=4.4cm, align=left, font=\footnotesize\sffamily, text=uefsInk] at (-2.25, 0.85) {%
    $\bullet$ Dinâmica regida por leis\\[-1pt]\hspace{0.55em}físico-biológicas\\[5pt]
    $\bullet$ Trabalho adapta-se aos\\[-1pt]\hspace{0.55em}ritmos naturais\\[5pt]
    $\bullet$ Baixa densidade técnica
  };
\end{scope}

% =========================================================
% PAINEL 2: Meio Técnico (Segunda Natureza)
% =========================================================
\begin{scope}[shift={(0, 1.05)}]
  \draw[draw=uefsGoldDark, fill=bgGold, rounded corners=6pt, line width=1.2pt] (-2.45, -1.75) rectangle (2.45, 1.85);
  \fill[uefsGoldDark, rounded corners=5pt] (-2.45, 1.05) rectangle (2.45, 1.85);
  \fill[uefsGoldDark] (-2.45, 1.05) rectangle (2.45, 1.35);
  \node[align=center, font=\small\bfseries\sffamily, text=white] at (0, 1.45)
    {2. Meio Técnico\\[-1pt]\footnotesize\normalfont(Segunda Natureza)};

  \node[anchor=north west, text width=4.4cm, align=left, font=\footnotesize\sffamily, text=uefsInk] at (-2.25, 0.85) {%
    $\bullet$ Natureza transformada pela\\[-1pt]\hspace{0.55em}indústria e capital\\[5pt]
    $\bullet$ Infraestruturas, barragens,\\[-1pt]\hspace{0.55em}cidades\\[5pt]
    $\bullet$ Racionalidade instrumental
  };
\end{scope}

% =========================================================
% PAINEL 3: Meio Técnico-Científico-Informacional
% =========================================================
\begin{scope}[shift={(5.8, 1.05)}]
  \draw[draw=uefsNavy, fill=bgNavy, rounded corners=6pt, line width=1.2pt] (-2.45, -1.75) rectangle (2.45, 1.85);
  \fill[uefsNavy, rounded corners=5pt] (-2.45, 1.05) rectangle (2.45, 1.85);
  \fill[uefsNavy] (-2.45, 1.05) rectangle (2.45, 1.35);
  \node[align=center, font=\small\bfseries\sffamily, text=white] at (0, 1.45)
    {3. Meio Técnico-Científico-\\[-1pt]Informacional};

  \node[anchor=north west, text width=4.4cm, align=left, font=\footnotesize\sffamily, text=uefsInk] at (-2.25, 0.85) {%
    $\bullet$ Ciência, tecnologia e dados\\[-1pt]\hspace{0.55em}comandam o território\\[5pt]
    $\bullet$ Redes globais, satélites,\\[-1pt]\hspace{0.55em}biotecnologia\\[5pt]
    $\bullet$ Conflitos socioambientais
  };
\end{scope}

% Setas horizontais de transição histórica
\draw[->, >=Stealth, line width=2pt, uefsForest] (-3.25, 1.05) -- (-2.55, 1.05);
\draw[->, >=Stealth, line width=2pt, uefsGoldDark] (2.55, 1.05) -- (3.25, 1.05);

% =========================================================
% BANNER / PÍLULA INFERIOR DE SÍNTESE
% =========================================================
\node[draw=uefsTeal, fill=bgTeal, rounded corners=12pt, line width=1.3pt,
      text width=16.0cm, minimum height=0.95cm, inner sep=8pt, align=center,
      font=\small\sffamily, text=uefsNavy] at (0, -1.65)
  {\textbf{\color{uefsTeal}Porto-Gonçalves \& Enrique Leff:} a crise ambiental nasce da apropriação desigual dos territórios.};

\end{tikzpicture}
\end{document}
